\section{Theoretischer Hintergrund} \label{sec:theorie}

Dieses Kapitel erläutert die theoretischen Grundlagen des Systems: RFID-Technologie, KI-gestützte Gesichtserkennung sowie datenschutzrechtliche Rahmenbedingungen. Ein kritischer Blick auf den Stand der Technik ergänzt die Darstellung.

\subsection{RFID-Technologie} \label{subsec:rfid}
Radio-Frequency Identification (RFID) ermöglicht die berührungslose Objektidentifikation via Radiowellen. Ein System besteht aus Transponder (Tag) und Lesegerät. Das Projekt nutzt passive, \textit{EM4100-kompatible RFID-Technologie} im Niederfrequenzbereich (125\,kHz), die Energie per induktiver Kopplung aus dem Feld des Lesers bezieht \cite{finkenzeller2015rfid, em4100_datasheet, iso18000_2}. Die 64-Bit-Datensätze der Read-only-Transponder enthalten eine 40-Bit-UID und werden Manchester-kodiert übertragen \cite{em4100_datasheet}. Da EM4100 keine Verschlüsselung bietet, wird die UID offen übertragen. Dies ist im Schulkontext akzeptabel, da die Karten kostengünstig und robust sind. Die Sicherheit resultiert primär aus der Zwei-Faktor-Verifikation (RFID + Gesichtserkennung), wodurch das Klonen einer Karte lediglich dem Kopieren einer Papier-Bescheinigung entspricht.

\subsection{KI-basierte Gesichtserkennung} \label{subsec:gesichtserkennung}
Moderne Gesichtserkennung nutzt Deep Learning, um Merkmale direkt aus Bilddaten zu extrahieren. Zur Detektion wird der effiziente \textit{SSD Mobilenet v1} (Single Shot MultiBox Detector) eingesetzt, der Echtzeit-Verarbeitung auf lokaler Hardware erlaubt \cite{vladmandic2023faceapi}. Nach der Detektion werden \textit{Face Embeddings} generiert – hochdimensionale Vektoren, die charakteristische Gesichtsmerkmale repräsentieren \cite{schroff2015facenet}. Identität wird über den euklidischen Abstand zwischen Embeddings und einen Schwellenwert definiert \cite{goodfellow2016deep}. TensorFlow.js ermöglicht die Ausführung dieser Modelle in Node.js ohne Cloud-Anbindung \cite{tensorflow2023}.

\subsection{Rechtlicher Rahmen und Datenschutz (DSGVO)} \label{subsec:datenschutz}
Biometrische Identifikation in Schulen unterliegt der DSGVO sowie dem Bayerischen Datenschutzgesetz (BayDSG) \cite{dsgvo2016, baydsg}. Biometrische Daten sind nach Art. 9 Abs. 1 DSGVO besonders geschützt. Die Verarbeitung erfordert eine ausdrückliche Einwilligung (Art. 9 Abs. 2 lit. a) oder ein erhebliches öffentliches Interesse (Art. 9 Abs. 2 lit. g).

Die Verhältnismäßigkeit wird durch das Zwei-Faktor-Prinzip gewahrt: RFID dient als Primärmerkmal, die Biometrie lediglich zur optionalen Verifikation. Gemäß \textit{Privacy by Design} (Art. 25 DSGVO) erfolgt die Datenverarbeitung rein lokal (On-Premise). Biometrische Merkmale verlassen das Schulnetzwerk nicht, was die informationelle Selbstbestimmung schützt.

\subsection{Ethische Betrachtung und organisatorische Auswirkungen} \label{subsec:ethik}
Die Einführung biometrischer Systeme erfordert eine Abwägung zwischen Sicherheit und Überwachungsdruck. Ein transparentes \textit{Opt-In-Verfahren} stellt sicher, dass Nutzer frei zwischen biometrischer Zusatzfunktion und rein kartenbasiertem System wählen können. Als Assistenzsystem für das Personal wahrt die Lösung den \enquote{Human-in-the-Loop}-Grundsatz der KI-Ethik-Leitlinien \cite{aiethics2019}.

Organisatorisch fallen initialer Aufwand für Schulungen und Aufklärung an. Langfristig reduziert die Automatisierung jedoch Kosten und manuellen Aufwand. Die Akzeptanz hängt vom Sicherheitsgewinn und der Zeitersparnis bei hoher Usability ab.

\subsection{Stand der Technik und existierende Systeme} \label{subsec:standdertechnik}
Gängige Verfahren reichen von manuellen Listen über analoge Ausweise bis zu Cloud-Lösungen. Manuelle Kontrollen sind fehleranfällig und kaum skalierbar. Analoge oder reine Kartensysteme bieten keine visuelle Inhaberverifizierung und sind bei Verlust unsicher. Zudem sind viele Chipsysteme (z. B. Mifare Classic) anfällig für Kloning-Angriffe.

Das entwickelte System kombiniert die Effizienz automatisierter Validierung mit lokaler biometrischer Verifizierung. Durch moderne Webtechnologien wie Express.js \cite{expressjs2023} und React Native \cite{reactnative2023} erreicht es eine Skalierbarkeit, die sonst proprietären Enterprise-Lösungen vorbehalten ist, während es die strengen Datenschutzvorgaben deutscher Schulen erfüllt.
